\begin{document}

%% Project-specific Preambles
\settoggle{includelogo}{true}

\newcommand{\thesubject}{Begabtenförderung}
\newcommand{\thetopic}{Risiko simulieren}
\newcommand{\thelogo}{Logos/FreeFlower.pdf}
\newcommand{\logosize}{.5\textwidth}
\newcommand{\figdir}{FF/Risiko_simulieren/Figures}
% Hilfestellung bleibt auch in der Version ohne Lösungen sichtbar
\newcommand{\hilfe}[1]{\par\medskip\noindent\textbf{Hilfestellung:} #1}
% Einstiegsfrage und Arbeitsblock je Doppellektion
\newcommand{\einstieg}[1]{\begin{tcolorboxTemplate}[title=Einstiegsfrage, colback=Goldenrod!10, colframe=Goldenrod!80!black, colbacktitle=Goldenrod!80!black]#1\end{tcolorboxTemplate}}
\newcommand{\arbeitsblock}[1]{\begin{myremark}\textbf{Arbeitsblock:} #1\end{myremark}}

\lhead[ \leftmark ]{\thetopic}
\rhead[\thesubject]{\textcopyleft{} \href{https://github.com/CyrilBlum/FreeFlower}{FreeFlower}, \the\year}

\settoggle{emoji}{false}
\def\arrayB{"grinning-face","bar-chart","laptop","technologist","robot","woman-student","hundred-points"}
\maketitlepage{\thetopic}{Tentativer Kursentwurf}
\tableofcontents
\newpage

\chapter{Kursidee und Ablauf}

\section{Worum es geht}
Die SchülerInnen lernen, wie sich Anlageformen unterscheiden, und simulieren selbst, was Unsicherheit über viele Jahre bewirkt. Jede Doppellektion beginnt mit einer Frage, auf die sie noch keine Antwort haben. Danach arbeiten sie selbstständig in vorbereiteten Notebooks. Am Schluss modellieren sie eine eigene reale Entscheidung und zeigen, wo ihr Modell versagt.

Der Kurs ist ein Spezialangebot der Informatik und wird von einer Lehrperson geleitet. Die Fragen stammen aus Wirtschaft und Finanzen, bearbeitet werden sie mit Simulationen in Python. Die nötigen Begriffe aus Wirtschaft und Recht führt der Kurs selbst ein, Vorwissen aus diesem Fach wird nicht vorausgesetzt. Der Kurs gibt keine Anlageempfehlungen.

\section{Umfang}
Der Kurs ist ein Förderangebot im Sinn des Konzepts Begabtenförderung. Er ergänzt den Regelunterricht und umfasst ein Quartal.

\begin{center}
	\small
	\begin{tabular}{@{}p{.3\textwidth}p{.6\textwidth}@{}}
		\midrule
		Präsenz & 8 Doppellektionen zu 90 Minuten, also 16 Lektionen \\
		Eigenarbeit & rund 8 Stunden für das Projekt \\
		Abschluss & Eigene Simulation mit Präsentation \\
		\midrule
	\end{tabular}
\end{center}

\section{Unterrichtsform}
Pro Doppellektion gibt es nur einen kurzen Impuls. Die SchülerInnen arbeiten zu zweit in einem Notebook, das Anleitung, Hilfestellungen und Selbstkontrollen enthält. Wer schneller ist, wählt die Challenges. Die detaillierte Planung steht im Anhang.

\begin{center}
	\small
	\begin{tabular}{@{}lrr@{}}
		\midrule
		\textbf{Phase} & \textbf{Minuten im ganzen Kurs} & \textbf{Anteil} \\
		\midrule
		Einstiegsfrage und Impuls & 105 & 15 \% \\
		Selbstständige Arbeit und Präsentationen & 510 & 71 \% \\
		Austausch in der Gruppe & 105 & 15 \% \\
		\midrule
	\end{tabular}
\end{center}

\section{Material}
Die Notebooks liegen unter \texttt{FF/Risiko\_simulieren/Notebooks}, die Lösungen im Unterordner \texttt{Loesungen}. Die Notebooks 1 bis 7 gehören zu den Doppellektionen 1 bis 7. Notebook 8 ist die Projektvorlage.

\chapter{Zufall simulieren}

\einstieg{Wie gross ist die Wahrscheinlichkeit, dass Sie mit 80 kein Geld mehr haben, und wer ist dann schuld? Heute können Sie das nicht beantworten. In der letzten Doppellektion sollten Sie zeigen können, wie man es abschätzt.\par\medskip Für heute: Ein Casino verliert an rund jedem dritten Abend gegen eine Person, die stur auf Rot setzt. Weshalb gehört es trotzdem nie dem Pech?}

\arbeitsblock{Notebook \texttt{01\_Zufall.ipynb}, 60 Minuten zu zweit. Die Selbstkontrollen im Notebook zeigen, ob Ihr Ergebnis stimmt.}

\section{Spiel zum Einstieg}
Jede Person startet mit 100 Punkten und setzt vor jedem Wurf einen selbst gewählten Teil davon. Bei einer geraden Zahl gewinnt sie den Einsatz, bei 1 oder 3 verliert sie ihn, bei 5 verliert sie das Doppelte. Nach zehn Runden vergleichen wir die Kontostände und besprechen, wer wie entschieden hat.

\section{Vom einzelnen Wurf zum Erwartungswert}
Ein einzelner Wurf sagt wenig. Über viele Wiederholungen nähert sich der Durchschnitt aber dem \tib{Erwartungswert}, also der Summe aller möglichen Werte, gewichtet mit ihren Wahrscheinlichkeiten. Eine \tib{Monte-Carlo-Simulation} nutzt genau das: Der Computer spielt eine Situation sehr oft durch, und wir werten die Ergebnisse aus.

\section{Das erste Notebook}
Ein Jupyter Notebook besteht aus Zellen mit Text oder Code. Sie führen eine Zelle aus und sehen das Ergebnis direkt darunter. Das eignet sich für Simulationen, weil Sie einen Parameter ändern und das Ergebnis sofort vergleichen können.

\begin{lstlisting}[style=Python]
import numpy as np

rng = np.random.default_rng(1)              # fester Startwert
wurf = rng.integers(1, 7, size=100_000)     # 100'000 Würfelwürfe
print(wurf.mean())                          # 3.49739
\end{lstlisting}

\begin{myexercise}[Roulette]
	Beim europäischen Roulette gibt es 18 rote, 18 schwarze und eine grüne Zahl. Sie setzen 1~CHF auf Rot und erhalten bei Rot Ihren Einsatz plus 1~CHF zurück.
	\begin{aenum}
		\item Berechnen Sie den Erwartungswert Ihres Gewinns pro Spiel.
		\item Simulieren Sie 1000 Abende mit je 100 Spielen. Wie viele Abende enden im Plus?
	\end{aenum}
	\begin{myanswer}
		\begin{aenum}
			\item $\frac{18}{37}\cdot 1 + \frac{19}{37}\cdot(-1) = -\frac{1}{37} \approx -0.027$~CHF.
			\item Exakt sind es 35.5~\% der Abende (Binomialverteilung, mindestens 51 Treffer). Jedes einzelne Spiel kostet im Durchschnitt trotzdem Geld. Summiert über 50'000 Spiele eines Abends im Casino verschwindet das Glück der Einzelnen.
		\end{aenum}
	\end{myanswer}
\end{myexercise}

\begin{mychallenge}[Wie genau ist eine Simulation?]
	Wie viele Würfelwürfe brauchen Sie, damit der Durchschnitt mit 95~\% Sicherheit höchstens 0.01 vom Erwartungswert 3.5 abweicht?
	\hilfe{Der Durchschnitt von $n$ Werten streut mit dem Standardfehler $\sigma/\sqrt{n}$. Die 95-\%-Grenze liegt bei 1.96 Standardfehlern. Berechnen Sie zuerst $\sigma^2$ eines Würfels aus $E[X^2] - E[X]^2$.}
	\begin{myanswer}
		$\sigma^2 = \frac{91}{6} - 3.5^2 = \frac{35}{12}$. Aus $1.96\,\sigma/\sqrt{n} = 0.01$ folgen rund 112'000 Würfe. Für jede zusätzliche Dezimalstelle braucht es hundertmal mehr Simulationen.
	\end{myanswer}
\end{mychallenge}

\chapter{Anlageformen}\label{chap:anlageformen}

\einstieg{Was genau besitzen Sie, wenn Ihr Geld \enquote{auf der Bank} liegt?}

\arbeitsblock{In Doppellektion 2 erstellen Sie einen Steckbrief und arbeiten mit Notebook 2, in Doppellektion 3 mit Notebook 3.}

\section{Was man besitzt und wer das Risiko trägt}
Wer Geld anlegt, stellt es jemandem zur Verfügung und erhält dafür eine unsichere oder eine vereinbarte Gegenleistung.

\begin{center}
	\small
	\begin{tabular}{@{}p{.18\textwidth}p{.27\textwidth}p{.22\textwidth}p{.22\textwidth}@{}}
		\midrule
		\textbf{Anlageform} & \textbf{Was man besitzt} & \textbf{Ertrag} & \textbf{Hauptrisiko} \\
		\midrule
		Sparkonto & Forderung gegenüber der Bank & Zins & Teuerung, Ausfall der Bank \\
		Obligation & Forderung gegenüber Staat oder Unternehmen & Zins, Rückzahlung am Ende & Ausfall, steigende Zinsen \\
		Aktie & Anteil an einem Unternehmen & Dividende, Kursgewinn & Unternehmen und Markt \\
		Fonds, \acrshort{ETF} & Anteil an einem Korb von Anlagen & wie der Korb & Markt, Kosten \\
		Immobilie & Grundstück und Gebäude & Miete, Wertsteigerung & Klumpenrisiko, schlechte Handelbarkeit \\
		Kryptowährung & Einheit in einem Netzwerk & nur Kursgewinn & starke Kursschwankungen, Betrug \\
		\midrule
	\end{tabular}
\end{center}

Ein \gls{ETF} ist ein Fonds, der an der Börse gehandelt wird und meist einen Index abbildet.

\section{Rendite und Volatilität}
Die \tib{Rendite} misst, um wie viel Prozent eine Anlage in einem Zeitraum gewachsen ist. Die \tib{Volatilität} misst, wie stark die Renditen schwanken, und wird meist als Standardabweichung angegeben. Sie ist das übliche Mass für Risiko, deckt aber nicht alles ab, wie wir in \cref{chap:modell} sehen.

\begin{myexercise}[Zinseszins]
	\begin{aenum}
		\item Wie viel werden 1000~CHF in 30 Jahren bei konstant 7~\%?
		\item Eine Anlage verliert zuerst 50~\% und gewinnt danach 50~\%. Wie viel Prozent liegt sie unter dem Startwert?
	\end{aenum}
	\begin{myanswer}
		\begin{aenum}
			\item $1000\cdot 1.07^{30} \approx 7612$~CHF.
			\item $0.5 \cdot 1.5 = 0.75$, also 25~\% unter dem Startwert.
		\end{aenum}
	\end{myanswer}
\end{myexercise}

\section{Nicht eine Zukunft, sondern viele}
\einstieg{Zwei Personen sparen 40 Jahre lang genau gleich viel in genau denselben Fonds. Im Modell hat die eine am Ende rund zwölfmal so viel wie die andere. War sie klüger?}

Eine feste Rendite ist unrealistisch. Wir ziehen deshalb für jedes Jahr eine zufällige Rendite und wiederholen das für viele mögliche Verläufe.

\begin{lstlisting}[style=Python]
jahre, pfade = 30, 10_000
rendite = rng.normal(0.06, 0.16, size=(pfade, jahre))   # eine Zeile pro Verlauf
vermoegen = 1000 * np.exp(rendite.sum(axis=1))

print(np.median(vermoegen))              # typischer Verlauf
print(np.percentile(vermoegen, [5, 95])) # schlechte und gute 5 %
print(np.mean(vermoegen < 1000))         # Anteil mit Verlust
\end{lstlisting}

Die Parameter in \cref{tab:anlagen} sind Annahmen für den Unterricht und keine Prognosen. Die Ergebnisse stammen aus je 100'000 simulierten Verläufen.

\begin{table}[H]
	\centering
	\small
	\caption{Aus 1000~CHF werden nach 30 Jahren im Modell (Beträge gerundet)}
	\label{tab:anlagen}
	\begin{tabular}{@{}lrrrrrr@{}}
		\midrule
		\textbf{Anlage} & \textbf{Wachstum} & \textbf{Vola.} & \textbf{Schlechte 5 \%} & \textbf{Median} & \textbf{Gute 5 \%} & \textbf{Verlust} \\
		\midrule
		Sparkonto & 0.5 \% & 0 \% & 1160 & 1160 & 1160 & 0 \% \\
		Obligationen & 2 \% & 4 \% & 1270 & 1820 & 2610 & 0.3 \% \\
		Aktienfonds & 6 \% & 16 \% & 1430 & 6030 & 25'400 & 2 \% \\
		Krypto (extrem) & 10 \% & 80 \% & 14 & 20'000 & 26.7 Mio. & 25 \% \\
		\midrule
	\end{tabular}
\end{table}

\begin{figure}[H]
	\centering
	\begin{tikzpicture}
		\begin{axis}[width=.8\textwidth, height=6cm, xlabel={Jahre}, ylabel={Vermögen in CHF}, ymode=log, legend pos=north west, grid=major, xmin=0, xmax=30]
			\addplot[Red, thick] table[x=jahr, y=p95] {\figdir/faecher_aktien.dat};
			\addlegendentry{Gute 5 \%}
			\addplot[Blue, very thick] table[x=jahr, y=p50] {\figdir/faecher_aktien.dat};
			\addlegendentry{Median}
			\addplot[Gray, thick] table[x=jahr, y=p05] {\figdir/faecher_aktien.dat};
			\addlegendentry{Schlechte 5 \%}
		\end{axis}
	\end{tikzpicture}
	\caption{Aktienfonds im Modell: Die Bandbreite wächst mit der Zeit (logarithmische Achse)}
\end{figure}

\begin{myexercise}[Eigene Parameter]
	Ändern Sie Wachstum und Volatilität. Was verändert den Median, was die Bandbreite? Begründen Sie, weshalb die Zeile zur Krypto-Anlage mit Vorsicht zu lesen ist.
	\begin{myanswer}
		Das Wachstum verschiebt den Median, die Volatilität vergrössert die Bandbreite. Die Krypto-Zeile beruht auf gewählten Parametern und auf der Annahme, dass die Anlage 30 Jahre lang existiert. Ein Totalverlust ist im Modell gar nicht vorgesehen.
	\end{myanswer}
\end{myexercise}

\begin{mychallenge}[Mittelwert gegen Median]
	Der Median im Aktienmodell liegt bei rund 6050~CHF. Berechnen Sie den Mittelwert aus einer Formel und prüfen Sie ihn in der Simulation. Wie viele Personen liegen unter dem Durchschnitt?
	\hilfe{Das Endvermögen ist lognormalverteilt. Für $X = e^{Y}$ mit $Y \sim N(\mu, s^2)$ gilt $E[X] = e^{\mu + s^2/2}$. Hier sind $\mu = 30 \cdot 0.06$ und $s^2 = 30 \cdot 0.16^2$. Für den Anteil genügt \texttt{np.mean(vermoegen < vermoegen.mean())}.}
	\begin{myanswer}
		$1000 \cdot e^{1.8 + 0.384} \approx 8880$~CHF, die Simulation liefert praktisch denselben Wert. Rund 67~\% der Personen liegen unter dem Durchschnitt. Der Median ist also die ehrlichere Zahl für den Einzelfall.
	\end{myanswer}
\end{mychallenge}

\begin{mychallenge}[Sparplan oder alles auf einmal?]
	Sie haben 12'000~CHF. Entweder investieren Sie alles sofort oder monatlich 100~CHF während 10 Jahren. Vergleichen Sie Median, schlechte 5~\% und Verlustwahrscheinlichkeit bei 6~\% Wachstum und 16~\% Volatilität pro Jahr.
	\hilfe{Rechnen Sie in Monaten mit $0.06/12$ und $0.16/\sqrt{12}$. Die Rate aus Monat $t$ wächst nur über die verbleibenden Monate. \texttt{np.cumsum(r[:, ::-1], axis=1)[:, ::-1]} liefert diese Restsummen.}
	\begin{myanswer}
		Einmalanlage: Median rund 21'900~CHF, schlechte 5~\% rund 9500~CHF, Verlust in 12~\% der Fälle. Sparplan: Median rund 16'700~CHF, schlechte 5~\% rund 10'200~CHF, Verlust in 14~\% der Fälle. Die Einmalanlage gewinnt in 84~\% der Verläufe, weil das Geld länger arbeitet. Der Sparplan senkt also nicht jedes Risiko.
	\end{myanswer}
\end{mychallenge}

\begin{myexcursion}
	In der Simulation addieren wir die Renditen und wenden danach die Exponentialfunktion an. Das entspricht stetig verzinsten Renditen und verhindert, dass ein Vermögen negativ wird.
\end{myexcursion}

\chapter{Der Durchschnitt täuscht}

\einstieg{In diesem Spiel wird der Durchschnitt aller Teilnehmenden in jeder Runde reicher. Weshalb sollten Sie trotzdem nicht mitspielen?}

\arbeitsblock{Notebook \texttt{04\_Durchschnitt.ipynb}, 60 Minuten zu zweit.}

\section{Ein Spiel mit positivem Durchschnitt}
Sie starten mit 100~CHF. In jeder Runde wird eine Münze geworfen. Bei Kopf wächst Ihr Vermögen um 50~\%, bei Zahl sinkt es um 40~\%. Würden Sie 20 Runden spielen?

Pro Runde beträgt der Erwartungswert des Faktors $0.5\cdot 1.5 + 0.5\cdot 0.6 = 1.05$. Das klingt nach 5~\% Gewinn pro Runde. Eine Runde Kopf und eine Runde Zahl ergeben zusammen aber $1.5\cdot 0.6 = 0.9$, also einen Verlust.

\begin{lstlisting}[style=Python]
faktor = rng.choice([1.5, 0.6], size=(100_000, 20))
ende = 100 * faktor.prod(axis=1)
print(ende.mean(), np.median(ende))    # 264.87... 34.86...
\end{lstlisting}

Mit 200'000 simulierten Personen erhalten wir: Der Durchschnitt liegt bei rund 264~CHF, der theoretische Wert bei 265.30~CHF. Die mittlere Person besitzt am Schluss aber nur 34.87~CHF. Rund 75~\% der Personen haben Geld verloren, und nur rund 13~\% liegen über dem Durchschnitt. Wenige sehr grosse Vermögen ziehen den Durchschnitt nach oben.

\begin{figure}[H]
	\centering
	\begin{tikzpicture}
		\begin{axis}[width=.8\textwidth, height=6cm, xlabel={Runde}, ylabel={Vermögen in CHF}, ymode=log, legend pos=north west, grid=major, xmin=0, xmax=20]
			\addplot[Red, very thick] table[x=runde, y=theorie] {\figdir/muenzspiel.dat};
			\addlegendentry{Durchschnitt aller Personen}
			\addplot[Blue, very thick] table[x=runde, y=median] {\figdir/muenzspiel.dat};
			\addlegendentry{Median}
		\end{axis}
	\end{tikzpicture}
	\caption{Durchschnitt und Median laufen auseinander (gerade Runden, logarithmische Achse)}
\end{figure}

\begin{myexercise}[Warum die Wachstumsrate zählt]
	\begin{aenum}
		\item Berechnen Sie die typische Wachstumsrate pro Runde.
		\item Was ändert sich für eine einzelne Person, wenn sie nur jeweils 40~\% ihres Vermögens einsetzt? Prüfen Sie es in der Simulation.
	\end{aenum}
	\begin{myanswer}
		\begin{aenum}
			\item $\sqrt{1.5\cdot 0.6} \approx 0.949$, also rund $-5.1$~\% pro Runde.
			\item Der Faktor ist nun $1.2$ oder $0.84$. Es gilt $\sqrt{1.2\cdot 0.84} \approx 1.004$, die typische Person wächst also leicht. Wer weniger einsetzt, verdient im Durchschnitt weniger und steht im Normalfall besser da.
		\end{aenum}
	\end{myanswer}
\end{myexercise}

\begin{mychallenge}[Der beste Einsatz]
	Sie setzen pro Runde den Anteil $f$ Ihres Vermögens ein. Bei welchem $f$ wächst das typische Vermögen am schnellsten, und ab welchem $f$ schrumpft es?
	\hilfe{Die typische Wachstumsrate ist $g(f) = \sqrt{(1+0.5f)(1-0.4f)}-1$. Bestimmen Sie das Maximum mit \texttt{np.linspace} und \texttt{argmax}. Wer Ableitungen kennt, leitet $\ln(1+0.5f) + \ln(1-0.4f)$ nach $f$ ab und setzt null.}
	\begin{myanswer}
		Das Maximum liegt bei $f = 0.25$ mit rund $+0.62$~\% pro Runde. Ab $f = 0.5$ schrumpft das Vermögen, denn dann gilt $(1+0.25)(1-0.2) = 1$. Das ist das Kelly-Kriterium: Vollen Einsatz bezahlt man mit dem sicheren Abstieg, obwohl jeder Einsatz einen positiven Erwartungswert hat.
	\end{myanswer}
\end{mychallenge}

\section{Reihenfolge der Renditen}
Wer sein Vermögen unberührt lässt, kann die Reihenfolge der Jahresrenditen ignorieren, denn eine Multiplikation lässt sich vertauschen. Sobald Geld entnommen wird, gilt das nicht mehr.

\begin{myexercise}[Rente aus dem Vermögen]
	Eine Person besitzt 100'000~CHF und entnimmt jedes Jahr zu Beginn 12'000~CHF. Die zehn Jahresrenditen lauten 20, 15, 10, 10, 5, 5, 0, $-5$, $-15$ und $-25$~\%. Simulieren Sie diese Reihenfolge und die umgekehrte.
	\begin{myanswer}
		Mit den guten Jahren am Anfang bleiben rund 21'600~CHF übrig. Mit den schlechten Jahren am Anfang reicht das Geld nicht, und der Rechnung fehlen am Schluss rund 75'100~CHF. Ohne Entnahme wäre der Endwert in beiden Fällen gleich.
	\end{myanswer}
\end{myexercise}

\chapter{Diversifikation}

\einstieg{Sie besitzen 500 verschiedene Aktien. Weshalb können Sie trotzdem an einem Tag einen grossen Teil verlieren?}

\arbeitsblock{Notebook \texttt{05\_Diversifikation.ipynb}, 35 Minuten. Danach 25 Minuten Projektidee festlegen.}

\section{Viele Körbe, wenig Schutz im Sturm}
Wer sein Geld auf viele Anlagen verteilt, gleicht einen Teil der Schwankungen aus. Wie viel das bringt, hängt davon ab, wie stark die Anlagen zusammenhängen. Dieses Mass heisst \tib{Korrelation}.

Für ein gleich gewichtetes Portfolio aus $n$ Aktien mit gleicher Volatilität $\sigma$ und gleicher Korrelation $\rho$ gilt
\[
	\sigma_{\text{Portfolio}} = \sigma\sqrt{\frac{1}{n} + \frac{n-1}{n}\rho}.
\]
Für grosse $n$ bleibt $\sigma\sqrt{\rho}$ übrig. Dieser Teil lässt sich nicht wegdiversifizieren.

\begin{figure}[H]
	\centering
	\begin{tikzpicture}
		\begin{axis}[width=.8\textwidth, height=6cm, xlabel={Anzahl Aktien}, ylabel={Volatilität in \%}, legend pos=north east, grid=major, xmin=1, xmax=50, ymin=0, ymax=32]
			\addplot[Blue, very thick] table[x=n, y=normal] {\figdir/diversifikation.dat};
			\addlegendentry{Korrelation 0.2 (ruhige Zeit)}
			\addplot[Red, very thick] table[x=n, y=krise] {\figdir/diversifikation.dat};
			\addlegendentry{Korrelation 0.7 (Krise)}
		\end{axis}
	\end{tikzpicture}
	\caption{Volatilität eines Portfolios bei $\sigma = 30$~\%}
\end{figure}

\begin{myexercise}[Formel und Simulation]
	Simulieren Sie 20 Aktien mit $\sigma = 30$~\% und vergleichen Sie mit der Formel. Erzeugen Sie die Korrelation mit einem gemeinsamen Marktfaktor.
	\begin{myanswer}
		Die Formel ergibt bei $\rho = 0.2$ rund 14.7~\% und bei $\rho = 0.7$ rund 25.4~\%. Die Simulation mit 200'000 Jahren liefert 14.69~\% und 25.40~\%. Die Grenzwerte liegen bei 13.4~\% und 25.1~\%. In einer Krise schützt die Verteilung auf 20 Aktien also deutlich weniger.
	\end{myanswer}
\end{myexercise}

\begin{myexercise}[Diskussion]
	Viele Korrelationen steigen in Krisen, weil viele AnlegerInnen gleichzeitig verkaufen. Weshalb ist das genau dann ein Problem, wenn man Diversifikation am meisten bräuchte?
	\begin{myanswer}
		Der Schutz wird in ruhigen Zeiten gemessen, gebraucht wird er aber in Krisen. Ein Modell mit konstanter Korrelation unterschätzt deshalb die Verluste im schlechten Fall.
	\end{myanswer}
\end{myexercise}

\begin{mychallenge}[Das risikoärmste Portfolio]
	Aktien haben $\sigma_1 = 16$~\%, Obligationen $\sigma_2 = 4$~\%, die Korrelation beträgt 0.1. Bei welchem Aktiengewicht $w$ ist die Volatilität am kleinsten? Weshalb ist die Antwort nicht $w = 0$?
	\hilfe{Verwenden Sie $\sigma_P^2 = w^2\sigma_1^2 + (1-w)^2\sigma_2^2 + 2w(1-w)\rho\sigma_1\sigma_2$ und testen Sie viele $w$. Prüfen Sie das Ergebnis mit zwei korrelierten Zufallsreihen.}
	\begin{myanswer}
		$w \approx 3.7$~\% ergibt eine Volatilität von rund 3.96~\%. Das liegt unter den 4~\% reiner Obligationen. Eine kleine Beimischung einer viel riskanteren Anlage senkt das Risiko also, weil die beiden nicht perfekt zusammenhängen.
	\end{myanswer}
\end{mychallenge}

\chapter{Versicherung im Alltag}\label{chap:franchise}

\einstieg{Wer eine hohe Franchise wählt, spart Prämie, weil er voraussichtlich gesund bleibt. Wer zahlt dann die Rechnung der Kranken, und ist das gerecht?}

\arbeitsblock{Notebook \texttt{06\_Franchise.ipynb}, 35 Minuten. Danach 25 Minuten Projektarbeit mit \texttt{08\_Projektvorlage.ipynb}.}

\section{Welche Franchise passt?}
In der obligatorischen Krankenversicherung zahlen Erwachsene ihre Kosten bis zur gewählten \tib{Franchise} selbst. Darüber übernehmen sie 10~\% als Selbstbehalt, höchstens 700~CHF pro Jahr. Eine höhere Franchise senkt die Prämie. Sie tauschen also sichere Kosten gegen unsichere.

Im Modell gelten fiktive Annahmen. Mit Franchise 300 beträgt die Jahresprämie 4200~CHF, und jeder zusätzliche Franken Franchise senkt sie um 70~Rappen. Die Behandlungskosten sind zufällig: Die gesunde Person hat einen Median von 600~CHF, die chronisch kranke einen von 6000~CHF. Die Simulation umfasst je eine Million Jahre.

\begin{figure}[H]
	\centering
	\begin{tikzpicture}
		\begin{axis}[width=.6\textwidth, height=6.4cm, xlabel={Franchise in CHF}, ylabel={Jahreskosten in CHF}, legend pos=outer north east, grid=major, xtick={300,500,1000,1500,2000,2500}, ymin=3400, ymax=6000, legend style={font=\small}]
			\addplot[Blue, very thick, mark=*] table[x=franchise, y=gesund] {\figdir/franchise.dat};
			\addlegendentry{Gesund, Durchschnitt}
			\addplot[Blue, dashed, mark=o] table[x=franchise, y=gesundp95] {\figdir/franchise.dat};
			\addlegendentry{Gesund, teure 5 \%}
			\addplot[Red, very thick, mark=*] table[x=franchise, y=chronisch] {\figdir/franchise.dat};
			\addlegendentry{Chronisch, Durchschnitt}
			\addplot[Red, dashed, mark=o] table[x=franchise, y=chronischp95] {\figdir/franchise.dat};
			\addlegendentry{Chronisch, teure 5 \%}
		\end{axis}
	\end{tikzpicture}
	\caption{Prämie plus Selbstbezahltes im Modell}
\end{figure}

Für die gesunde Person sinken die durchschnittlichen Kosten von 4559~CHF bei Franchise 300 auf 3651~CHF bei Franchise 2500. In 81~\% der Jahre ist Franchise 2500 günstiger. Für die chronisch kranke Person steigen sie von 5004~CHF auf 5426~CHF, Franchise 2500 lohnt sich nur in 9~\% der Jahre.

\begin{myexercise}[Schwelle]
	Ab welchen Behandlungskosten ist Franchise 300 günstiger als Franchise 2500?
	\begin{myanswer}
		Franchise 2500 spart 1540~CHF Prämie. Unter 2500~CHF Kosten zahlt man mit Franchise 2500 die Kosten $K$ selbst, mit Franchise 300 zahlt man $300 + 0.1(K-300) = 270 + 0.1K$. Aus $K - (270 + 0.1K) = 1540$ folgt $K \approx 2011$~CHF.
	\end{myanswer}
\end{myexercise}

\begin{myexercise}[Mit echten Zahlen]
	Suchen Sie im Prämienrechner des Bundes echte Prämien für Ihre Region und ein gewähltes Modell. Ersetzen Sie die fiktiven Werte. Was fehlt dem Modell trotzdem?
	\begin{myanswer}
		Zum Beispiel die Frage, ob eine Person hohe Kosten im Moment bezahlen kann. Dazu kommen Kosten, die über mehrere Jahre zusammenhängen, und die Tatsache, dass die eigene Gesundheit nicht einfach zufällig ist.
	\end{myanswer}
\end{myexercise}

\begin{mychallenge}[Wie gross muss eine Versicherung sein?]
	Jede Person hat pro Jahr mit 1~\% Wahrscheinlichkeit einen Schaden von 100'000~CHF. Welche Prämie deckt die Schäden in 99~\% der Jahre, wenn die Versicherung 100, 1000, 10'000 oder 100'000 Personen hat?
	\hilfe{Die Anzahl Schäden ist binomialverteilt, also verwenden Sie \texttt{rng.binomial} mit einer Million Wiederholungen. Nehmen Sie davon das 99-\%-Perzentil und teilen Sie die Schadensumme durch $n$. Der faire Beitrag liegt bei 1000~CHF.}
	\begin{myanswer}
		Die Prämie beträgt 4000, 1800, 1240 und 1074~CHF. Der Zuschlag schrumpft ungefähr mit $1/\sqrt{n}$. Das ist der Grund, weshalb Versicherungen gross sein wollen. Die Annahme unabhängiger Schäden fällt aber weg, sobald ein Hochwasser alle gleichzeitig trifft.
	\end{myanswer}
\end{mychallenge}

\chapter{Wenn das Modell versagt}\label{chap:modell}

\einstieg{Ein Modell sagt: Ein solcher Tag kommt alle 13'800 Jahre. Dann passiert er zweimal in zehn Jahren. Wer haftet?}

\arbeitsblock{Notebook \texttt{07\_Modell.ipynb}, 35 Minuten. Danach 25 Minuten Projektarbeit.}

\section{Dicke Ränder}
Bisher haben wir meist normalverteilte Renditen angenommen. Reale Kursbewegungen zeigen aber häufiger extreme Tage, als die Normalverteilung erwartet. Wir vergleichen zwei Modelle mit gleicher Tagesvolatilität von 1~\% über 2520 Handelstage, also rund zehn Jahre. Das zweite Modell verwendet eine t-Verteilung mit 3 Freiheitsgraden.

\begin{table}[H]
	\centering
	\small
	\caption{Durchschnitt aus 2000 simulierten Zehnjahres-Perioden}
	\label{tab:extrem}
	\begin{tabular}{@{}lrrrr@{}}
		\midrule
		\textbf{Modell} & \textbf{Tage $\leq -3$ \%} & \textbf{Tage $\leq -5$ \%} & \textbf{\acrshort{VaR} 99 \%} & \textbf{Schlechtester Tag} \\
		\midrule
		Normalverteilung & 3.5 & 0.0 & $-2.3$ \% & $-3.4$ \% \\
		t-Verteilung & 17.2 & 4.0 & $-2.6$ \% & $-9.0$ \% \\
		\midrule
	\end{tabular}
\end{table}

Der \gls{VaR} 99~\% gibt den Tagesverlust an, der nur an einem von 100 Tagen überschritten wird. Beide Modelle liegen dabei nahe beieinander, obwohl sich die Extremtage deutlich unterscheiden. Der \acrshort{VaR} sagt nichts darüber, wie schlimm die übrigen Tage ausfallen. Beim schlechtesten Tag ist der Median angegeben.

\begin{myexercise}[Fünf Standardabweichungen]
	Ein Verlust von 5~\% entspricht hier fünf Standardabweichungen. Die Normalverteilung ordnet dem eine Wahrscheinlichkeit von rund $2.9\cdot 10^{-7}$ zu. Wie oft käme ein solcher Tag bei 252 Handelstagen pro Jahr vor?
	\begin{myanswer}
		Etwa einmal in $1/(2.9\cdot 10^{-7}\cdot 252)$, also rund 13'800 Jahren. Wenn reale Märkte solche Tage innerhalb weniger Jahre zeigen, liegt das Problem beim Modell.
	\end{myanswer}
\end{myexercise}

\begin{mychallenge}[Schlimmer als der Value at Risk]
	Der Expected Shortfall 99~\% ist der Durchschnitt aller Tage, die schlechter sind als der \acrshort{VaR}. Berechnen Sie beide Zahlen für beide Modelle.
	\hilfe{\texttt{v = np.percentile(r, 1)} liefert den \acrshort{VaR}, \texttt{r[r <= v].mean()} den Expected Shortfall.}
	\begin{myanswer}
		Der \acrshort{VaR} liegt bei $-2.3$~\% und $-2.6$~\%, der Expected Shortfall bei $-2.7$~\% und $-4.0$~\%. Der Expected Shortfall trennt die Modelle deutlich besser, weil er in den Rand hineinschaut.
	\end{myanswer}
\end{mychallenge}

\begin{mychallenge}[Schulden als Verstärker]
	Ein Fonds hat eine Tagesvolatilität von 0.5~\%, eine Tagesdrift von 0.04~\% und arbeitet mit Hebel $h$. Wie oft verliert er innerhalb eines Jahres 80~\% seines Eigenkapitals?
	\hilfe{Die Tagesrendite des Eigenkapitals ist $h \cdot r$. Rechnen Sie mit Logarithmen, denn die lassen sich mit \texttt{np.cumsum} aufsummieren. Im Notebook steht der passende Ansatz. Vergleichen Sie die Summe mit \texttt{np.log(0.2)}.}
	\begin{myanswer}
		\begin{center}
			\small
			\begin{tabular}{@{}lrrr@{}}
				\midrule
				\textbf{Modell} & \textbf{Hebel 5} & \textbf{Hebel 10} & \textbf{Hebel 25} \\
				\midrule
				Normalverteilung & 0.0 \% & 0.5 \% & 31 \% \\
				t-Verteilung & 0.2 \% & 1.9 \% & 34 \% \\
				\midrule
			\end{tabular}
		\end{center}
		Bei Hebel 10 liegt das Risiko mit dicken Rändern rund viermal höher. Bei Hebel 25 reicht schon die normale Schwankung. Der Hebel wirkt stärker als die Wahl der Verteilung.
	\end{myanswer}
\end{mychallenge}

\section{Fallbeispiele}
Die SchülerInnen wählen in Gruppen einen Fall und klären drei Fragen: Was hat das Modell angenommen, was ist tatsächlich passiert, und wer trug den Schaden?

\begin{itemize}
	\item \textbf{Frankenschock, Januar 2015:} Die Schweizerische Nationalbank hebt den Mindestkurs zum Euro auf. Ein Kurs, der lange fast fest war, springt innert kurzer Zeit. Was misst eine Volatilität, die vor dem Ereignis geschätzt wurde?
	\item \textbf{Long-Term Capital Management, 1998:} Ein Hedgefonds mit renommierten Finanzökonomen arbeitet mit hohen Schulden und gerät nach der Russlandkrise in Schieflage. Banken organisieren die Rettung. Weshalb verstärken Schulden Modellfehler?
	\item \textbf{Finanzkrise 2008:} Hypothekenkredite galten durch Bündelung als gut gestreut. Dann fielen viele gleichzeitig aus. Das schliesst direkt an die Korrelation in Kapitel 5 an: Wie unabhängig waren die Risiken wirklich?
\end{itemize}

\begin{myexercise}[Modellkritik]
	Nennen Sie drei Annahmen aus Ihrer bisherigen Simulation, die in der Realität verletzt sein könnten, und testen Sie eine davon im Notebook.
	\begin{myanswer}
		Beispiele sind unabhängige Jahre, konstante Volatilität, normalverteilte Renditen, keine Kosten und keine Steuern, oder dass eine Anlage immer weiterbesteht.
	\end{myanswer}
\end{myexercise}

\chapter{Projekte und weitere Ideen}

\einstieg{Welche Entscheidung möchten Sie untersuchen, und wem würde ein Fehler in Ihrem Modell schaden?}

\section{Projektauftrag}
Sie wählen zu zweit eine reale Entscheidung unter Unsicherheit und bauen dazu eine eigene Simulation in \texttt{08\_Projektvorlage.ipynb}. Abzugeben sind ein lauffähiges Notebook, eine Grafik mit Ihrer Kernaussage und eine Präsentation von zehn Minuten. Die Präsentation beantwortet diese Fragen:
\begin{itemize}
	\item Welche Entscheidung untersuchen Sie, und für wen ist sie wichtig?
	\item Welche Annahmen trifft Ihr Modell, und woher stammen die Zahlen?
	\item Was zeigt die Simulation, auch im schlechten Fall?
	\item Bei welcher Änderung einer Annahme kippt Ihre Schlussfolgerung?
\end{itemize}

Die erwarteten Resultate unten stammen aus Referenzsimulationen mit den angegebenen Beispielparametern (\texttt{Code/referenzloesungen.py}). Eigene Parameter führen zu anderen Zahlen. Die Zahlen sind ein Vergleichswert, kein Ziel.

\newcommand{\projekt}[6]{%
	\subsection{#1}
	\begin{description}[leftmargin=!,labelwidth=2.6cm,topsep=2pt,itemsep=2pt]
		\item[Leitfrage] #2
		\item[Produkt] #3
		\item[Hilfestellung] #4
		\item[Erwartet] #5
		\item[Quellen] #6
	\end{description}}

\section{Projektideen}

\projekt{Altersvorsorge}
{Welche Aktienquote passt zu 40 Jahren Sparen, und wie schlimm ist der schlechte Fall?}
{Tabelle und Fächerdiagramm für Aktienquoten von 0 bis 100~\%.}
{Sparen Sie 6000~CHF pro Jahr. Mischen Sie pro Jahr \texttt{q * aktien + (1-q) * obligationen} und zinsen Sie jede Einzahlung über die Restjahre auf. Beginnen Sie mit den Parametern aus \cref{tab:anlagen}.}
{Median rund 372'000~CHF bei 0~\% Aktien und 1.08~Mio.~CHF bei 100~\%. Die schlechten 5~\% bleiben ab 25~\% Aktien fast gleich bei rund 340'000~CHF, die guten 5~\% steigen von rund 490'000~CHF auf 4~Mio.~CHF. Kritisch zu diskutieren ist die Annahme unabhängiger Jahre.}
{Bengen (1994), Determining Withdrawal Rates Using Historical Data, Journal of Financial Planning 7(4). Daten: \url{https://data.snb.ch/de}}

\projekt{Kelly-Kriterium}
{Welcher Einsatz maximiert das langfristige Wachstum, und weshalb setzen Profis weniger?}
{Diagramm mit dem Median und dem Risiko eines Einbruchs von 50~\% für verschiedene Einsätze.}
{Beginnen Sie mit Trefferquote 60~\%, Quote 1:1 und 1000 Runden. Rechnen Sie mit Logarithmen, sonst läuft die Zahl über. Die Kelly-Formel lautet hier $f^* = 2p - 1$. Testen Sie $f$ von 0.05 bis 0.5.}
{Das Maximum liegt bei $f = 0.2$, der Median vermehrt das Vermögen dort rund $5.6 \cdot 10^{8}$-fach. Bei $f = 0.4$ endet die mittlere Person unter dem Startwert. Selbst bei $f = 0.2$ brechen 45~\% der Verläufe zwischendurch um die Hälfte ein, bei $f = 0.1$ nur 10~\%. Das erklärt den halben Kelly.}
{Kelly (1956), A New Interpretation of Information Rate, Bell System Technical Journal 35(4). Peters (2019), The ergodicity problem in economics, Nature Physics 15.}

\projekt{Ruin des Spielers}
{Wie lange überlebt man mit begrenztem Kapital gegen einen kleinen Vorteil des Casinos?}
{Diagramm der Erfolgswahrscheinlichkeit gegen die Höhe des Kapitals, dazu die Formel und der Vergleich mit der Simulation.}
{Start $a$, Ziel $2a$, Einsatz 1 auf Rot. Die Formel lautet $P = (1-q^a)/(1-q^{2a})$ mit $q = \frac{19}{18}$. Simulieren Sie mit einer Schleife, die nur noch aktive Personen weiterspielen lässt.}
{Bei $a = 10$ verdoppeln rund 37~\%, bei $a = 50$ rund 6~\%, bei $a = 100$ nur noch 0.4~\%. Bei einem fairen Spiel wären es immer 50~\%. Je mehr Spiele es bis zum Ziel braucht, desto sicherer verliert man.}
{Feller, An Introduction to Probability Theory and Its Applications, Band 1, Kapitel XIV.}

\projekt{Bank Run}
{Ab wann kippt eine Bank, wenn KundInnen auf das Verhalten anderer reagieren?}
{Agentenmodell mit Diagramm des Endabzugs gegen den Anfangsschock.}
{Jede Person hat eine Schwelle: Sie hebt ab, sobald dieser Anteil abgehoben hat. Ziehen Sie die Schwellen \texttt{sw} aus einer Normalverteilung mit Mittelwert 0.4 und Streuung 0.2. Aktualisieren Sie wiederholt den Anteil mit \texttt{max(schock, np.mean(sw <= anteil))}, bis sich nichts mehr ändert. Variieren Sie Schock und Streuung.}
{Bei einer Streuung von 0.2 bleibt der Abzug bei einem Schock bis 20~\% stehen. Bei 30~\% steigt er im Mittel auf rund 78~\%, bei 35~\% auf fast 100~\%. Bei einer Streuung von 0.25 kippt es schon bei 5~\% Schock oft. Kleine Unterschiede in der Verteilung ändern also das ganze Ergebnis.}
{Diamond, Dybvig (1983), Bank Runs, Deposit Insurance, and Liquidity, Journal of Political Economy 91(3). Granovetter (1978), Threshold Models of Collective Behavior, American Journal of Sociology 83(6).}

\projekt{Gebäudeversicherung}
{Wie viel Reserve braucht eine Versicherung für seltene Grossschäden?}
{Histogramm der Jahresschäden mit Prämie und Reserve.}
{Modellieren Sie 100'000 Gebäude. Die Anzahl Schäden pro Jahr ist Poisson-verteilt mit Mittelwert 200, die Schadenhöhe Pareto-verteilt: \texttt{(rng.pareto(1.5, k) + 1) * 20\_000}. Die Reserve deckt das 99.5-\%-Perzentil.}
{Durchschnittlich rund 12~Mio.~CHF Schäden pro Jahr, also rund 120~CHF Prämie pro Gebäude. Das 99.5-\%-Perzentil liegt bei rund 35~Mio.~CHF, die Reserve also bei rund 23~Mio.~CHF. Wenige Grossschäden bestimmen die Reserve.}
{Klugman, Panjer, Willmot, Loss Models. Daten: Sigma-Studien des Swiss Re Institute, online verfügbar.}

\projekt{Klimarisiko}
{Wie verändern häufigere Extremereignisse Prämie und Reserve?}
{Vergleich heute und Zukunft in einer Tabelle, mit einer Begründung der gewählten Häufigkeit.}
{Bauen Sie auf dem Modell der Gebäudeversicherung auf und erhöhen Sie die Häufigkeit um 50~\%. Belegen Sie Ihre Annahme mit einer Quelle und testen Sie zusätzlich höhere Schäden.}
{Die Prämie steigt im Modell um 50~\% auf rund 180~CHF, die Reserve nur um rund ein Drittel auf rund 31~Mio.~CHF. Der Vergleich macht sichtbar, dass die Reserve nicht proportional wächst.}
{IPCC, Sixth Assessment Report, Working Group I, Kapitel 11 (Weather and Climate Extreme Events).}

\projekt{Lawinenrisiko}
{Wie addieren sich kleine Einzelrisiken über viele Skitouren?}
{Diagramm des Lebenszeitrisikos gegen die Anzahl Touren, mit und ohne risikoreiche Tage.}
{Verwenden Sie $1-(1-p)^n$. Die Beispielannahme $p = 1/10'000$ pro Tour ist nicht empirisch. Teilen Sie die Touren anschliessend in normale Tage und Tage mit dreifachem Risiko auf.}
{Bei 10 Touren pro Jahr während 50 Jahren beträgt das Risiko rund 4.9~\%. Sind 20~\% der Tage dreimal so gefährlich, steigt es auf rund 6.8~\%. Das Resultat hängt am gewählten $p$, deshalb gehört eine Begründung dazu.}
{WSL-Institut für Schnee- und Lawinenforschung SLF: \url{https://www.slf.ch/de/lawinen/}. Munter, 3x3 Lawinen.}

\projekt{Startup-Investments}
{Weshalb lebt ein Portfolio oft von sehr wenigen Treffern?}
{Diagramm der Verteilung der Portfoliorendite und Anteil des besten Treffers.}
{30 Investments zu je 1~CHF, 65~\% Totalverlust, sonst Vielfaches \texttt{rng.pareto(1.3, k) + 1}. Sortieren Sie die Ergebnisse pro Portfolio mit \texttt{np.sort}.}
{Das Vielfache beträgt im Durchschnitt rund 1.5, im Median aber nur rund 1.0, und rund 52~\% der Portfolios verlieren Geld. Der beste Treffer bringt im Median rund 30~\% des Ertrags, die besten drei rund 59~\%.}
{Kerr, Nanda, Rhodes-Kropf (2014), Entrepreneurship as Experimentation, Journal of Economic Perspectives 28(3).}

\projekt{Sportwetten}
{Lässt sich aus den Quoten die Marge des Anbieters berechnen?}
{Berechnung der Marge über eine ganze Saison und Simulation einer Wettstrategie.}
{Für Quoten 2.10, 3.40 und 3.60 ist $\sum 1/q$ grösser als 1. Der Überschuss ist die Marge. Normieren Sie die Kehrwerte, um faire Wahrscheinlichkeiten zu schätzen. Laden Sie danach eine Saison als CSV-Datei.}
{Im Beispiel beträgt die Summe 1.048, die Auszahlungsquote also rund 95.4~\%. Wer 200-mal 1~CHF setzt, verliert im Mittel rund 9~CHF, liegt aber in rund 25~\% der Fälle im Plus. Das erklärt, weshalb Wetten sich wie Können anfühlt.}
{Daten: \url{https://www.football-data.co.uk/data.php}}

\projekt{Verlustaversion}
{Wie entscheiden Menschen im Vergleich zum Erwartungswert?}
{Eigene Umfrage mit Auswertung und Schätzung des Verlustfaktors.}
{Fragen Sie: Münzwurf, Sie verlieren 100~CHF oder gewinnen $X$. Ab welchem $X$ spielen Sie? Der Quotient $X/100$ schätzt den Verlustfaktor $\lambda$. Halten Sie die Umfrage anonym und ohne Personendaten, und vergleichen Sie kleine mit grossen Beträgen.}
{Ein rein am Erwartungswert orientierter Mensch spielt knapp über $X = 100$. Tversky und Kahneman schätzen $\lambda$ auf rund 2.25, Ihre Umfrage dürfte ebenfalls deutlich über 1 liegen. Die Stichprobe ist klein und nicht repräsentativ.}
{Kahneman, Tversky (1979), Prospect Theory, Econometrica 47(2). Tversky, Kahneman (1992), Advances in Prospect Theory, Journal of Risk and Uncertainty 5(4).}

\projekt{Bootstrapping}
{Was ändert sich, wenn wir statt der Normalverteilung echte Renditen ziehen?}
{Vergleich von drei Modellen: Normalverteilung, Einzelziehung und Blockziehung.}
{Laden Sie Monatsrenditen des Marktes. Einzelziehung: \texttt{x[rng.integers(0, n, n)]}. Blockziehung: Ziehen Sie zusammenhängende Blöcke von 24 Monaten. Messen Sie \texttt{np.corrcoef(r[:-1]**2, r[1:]**2)}.}
{In einer Testreihe mit ruhigen und unruhigen Phasen beträgt diese Korrelation im Original rund 0.15, nach Einzelziehung rund 0, nach Blockziehung rund 0.11. Einzelziehung erhält die Verteilung und zerstört die Abfolge. Echte Daten dürften ähnlich reagieren, das ist zu prüfen.}
{Efron (1979), Bootstrap Methods, Annals of Statistics 7(1). Daten: \url{https://mba.tuck.dartmouth.edu/pages/faculty/ken.french/data_library.html}}

\projekt{Optionspreise}
{Wie viel ist das Recht wert, eine Aktie später zu einem festen Preis zu kaufen?}
{Monte-Carlo-Preis mit Konfidenzintervall und Vergleich mit der Black-Scholes-Formel.}
{Beispiel: Kurs 100, Ausübungspreis 100, Zins 1~\%, Volatilität 20~\%, ein Jahr. Simulieren Sie den Endkurs mit $S_T = S e^{(r-\sigma^2/2)T + \sigma\sqrt{T}Z}$, berechnen Sie $\max(S_T - K, 0)$ und diskontieren Sie mit $e^{-rT}$.}
{Black-Scholes ergibt 8.43. Die Simulation liegt bei 1000 Läufen noch rund $\pm 0.8$ daneben, bei 10~Mio. Läufen rund $\pm 0.01$. Das bestätigt die Regel aus Kapitel 2 zur Genauigkeit.}
{Black, Scholes (1973), The Pricing of Options and Corporate Liabilities, Journal of Political Economy 81(3). Hull, Options, Futures, and Other Derivatives.}

\section{Weitere Ideen für den Kurs}
\begin{itemize}
	\item \textbf{Klassenbörse:} Die Gruppen handeln über mehrere Wochen in einem simulierten Markt, in dem die Lehrperson Nachrichten einstreut. Danach vergleichen wir Stil und Ergebnis mit reinem Zufall.
	\item \textbf{Zufall gegen Auswahl:} Wir vergleichen zufällig gezogene Portfolios mit einem bewussten Auswahlverfahren und fragen, wie viel Erfolg Können ist.
	\item \textbf{Gast aus der Praxis:} Eine Fachperson aus dem Risikomanagement einer Versicherung oder Bank erzählt, wie Modelle dort eingesetzt und kontrolliert werden.
	\item \textbf{Besuch an einer Hochschule:} Ein Fachbereich für Finanzmathematik oder Versicherungsmathematik stellt ein Thema vor. Das passt zur Vernetzung im Konzept Begabtenförderung.
	\item \textbf{KI als Prüfobjekt:} Die SchülerInnen lassen ein Sprachmodell eine Simulation schreiben und suchen darin die versteckten Annahmen und Fehler.
	\item \textbf{Wettbewerbe:} Die Wirtschafts-Olympiade kann als Anschluss für besonders interessierte SchülerInnen dienen.
\end{itemize}

\newpage
\appendix
\chapter{Lektionsplanung}

Jede Doppellektion dauert 90 Minuten. Die Doppellektionen 1 bis 7 folgen demselben Ablauf: 10 Minuten Einstiegsfrage, 5 Minuten Impuls, 60 Minuten Arbeitsblock, 15 Minuten Austausch. Die Doppellektion 8 besteht aus Präsentationen.

% Argumente: Nummer, frei, Thema, Einstiegsfrage, Impuls, Arbeitsblock, Austausch
\newcommand{\lplan}[7]{%
	\subsection*{DL #1: #3}
	\begin{center}
		\small
		\begin{tabular}{@{}p{.07\textwidth}p{.19\textwidth}p{.65\textwidth}@{}}
			\midrule
			10' & Einstiegsfrage & #4 \\
			5' & Impuls & #5 \\
			60' & Arbeitsblock & #6 \\
			15' & Austausch & #7 \\
			\midrule
		\end{tabular}
	\end{center}}

\lplan{1}{}{Zufall simulieren}
{Geld mit 80 und Casino: Antworten sammeln, keine Auflösung. Danach Würfelspiel im Raum.}
{Notebook öffnen, Zelle ausführen, Selbstkontrolle erklären.}
{Zu zweit \texttt{01\_Zufall.ipynb}: Roulette, Casino-Bilanz, Challenge Genauigkeit.}
{Weshalb verliert das Casino nie? Was bedeutet Durchschnitt für den Einzelfall?}

\lplan{2}{}{Anlageformen}
{Was besitzen Sie, wenn das Geld auf der Bank liegt? Kurze schriftliche Antworten.}
{Forderung, Beteiligung, Risiko.}
{Steckbrief einer Anlageform zu zweit (30'), dann \texttt{02\_Anlageformen.ipynb} (30').}
{Galerierundgang: Wer verliert in welcher Anlage zuerst?}

\lplan{3}{}{30 Jahre in die Zukunft}
{Zwei Personen mit gleichem Plan, zwölffacher Unterschied: War die eine klüger?}
{Was ein Modellparameter ist und weshalb er keine Prognose ist.}
{Zu zweit \texttt{03\_30\_Jahre.ipynb}: Verläufe, Perzentile, Challenges Median und Sparplan.}
{Wie würden Sie das Ergebnis einer Person ohne Mathematik erklären?}

\lplan{4}{}{Der Durchschnitt täuscht}
{Münzspiel mit der ganzen Gruppe: Wer spielt mit und weshalb?}
{Mittelwert und Median an einem Beispiel.}
{Zu zweit \texttt{04\_Durchschnitt.ipynb}: Spiel, bester Einsatz, Reihenfolge der Renditen.}
{Lässt sich die Frage \enquote{Geld mit 80} nun beantworten? Was fehlt noch?}

\lplan{5}{}{Diversifikation}
{500 Aktien und trotzdem grosser Verlust: Wie kann das sein?}
{Korrelation an einem Bild.}
{\texttt{05\_Diversifikation.ipynb} (35'), danach Projektidee wählen und zu zweit festhalten (25').}
{Projektideen in je einem Satz vorstellen.}

\subsection*{Zwischen DL 5 und DL 6}
Die Projektgruppen geben ihre Leitfrage ab.

\lplan{6}{}{Versicherung im Alltag}
{Wer zahlt die Rechnung der Kranken? Kurze Standpunkte, auch rechtlich.}
{Franchise, Selbstbehalt, Solidarität.}
{\texttt{06\_Franchise.ipynb} (35'), danach Projektarbeit in \texttt{08\_Projektvorlage.ipynb} (25').}
{Ist das Modellergebnis gerecht? Was sagt das Modell nicht?}

\lplan{7}{}{Wenn das Modell versagt}
{Alle 13'800 Jahre, und dann zweimal in zehn Jahren: Wer haftet?}
{Dicke Ränder an einer Grafik.}
{\texttt{07\_Modell.ipynb} (35'), danach Fallbeispiel in Gruppen oder Projektarbeit (25').}
{Kurzer Fallbericht: Annahme, Ereignis, Geschädigte. Auftrag für die Eigenarbeit.}

\subsection*{Zwischen DL 7 und DL 8}
Eigenarbeit am Projekt, rund 8 Stunden.

\subsection*{DL 8: Präsentationen}
\begin{center}
	\small
	\begin{tabular}{@{}p{.07\textwidth}p{.19\textwidth}p{.65\textwidth}@{}}
		\midrule
		70' & Präsentationen & Bis zu sieben Gruppen mit je 10 Minuten inklusive Rückfrage \\
		20' & Abschluss & Rückkehr zur Frage \enquote{Geld mit 80}, Rückmeldung zum Kurs \\
		\midrule
	\end{tabular}
\end{center}

\newpage
\printglossaries

\end{document}
